% Terminal note of the homily: everything the speech itself must not carry ---
% the audience, the spoken word count and the pace it implies, the relation to
% the three interpretations of the expansive study, the route by which each
% quoted or reported text reaches the page, and the exact loci --- followed by
% the References. Nothing here is read aloud.

\section*{Note on Sources and Delivery}

{\small

\textbf{Audience and occasion.} An adult parish assembly, at the Mass of the
First Sunday of Advent, Year~B: \latin{Dominica I Adventus} of the \work{Roman
Missal}, Third Edition, for the dioceses of the United States, with no.~2 of
the \work{Lectionary for Mass}, for the occurrence of 29~November 2026. The
speech holds at every lawful celebration of the formulary: it names the
Communion verse only as the one the Missal gives, assumes no sung or recited
form of the psalm, the Alleluia verse or either antiphon, and does not
mention the solemn blessing, which is optional.

\textbf{Length and pace.} The spoken body runs to 1,398 words, counted as
whitespace-separated words once comments and markup are removed: about ten
and three-quarter minutes at an unhurried 130 words a minute, a little over
eleven at 125, and a little over eleven and a half at 120. That is
arithmetic on the count; no delivery has been timed. The prose was read
through in full in its final form for sense, sentence length and ease of
speech, silently, not aloud.

\textbf{Relation to the reviewed interpretations.} The argument preached is
the third of the expansive study's three interpretations, \textit{The Clay
and the Faithful Potter}, joined with the second, \textit{The Watch at the
Door}, as the study's comparison joins them: the third keeps the second from
sounding as though the servants secured their own readiness, and the second
keeps the third from passivity. From the second the speech takes the Gospel's
command to all and Augustine's reason for it, and Bede's servants given
authority over every work and his door of each heart. From the third it takes
the confession of Isaiah 63--64, Jerome on God's patience and on the clay,
Augustine on the psalm's refrain and on the Communion verse, the Collect's
petition for a will, and Chrysostom on Paul's thanksgiving and ``God is
faithful''. From the first, \textit{Rend the Heavens and Come Down}, it takes
only what all three share: the Preface's two comings. The past tense
\latin{Descendisti}, on which the first interpretation turns, is not
preached, because the English proclaimed in the United States does not read
it; the comparison finds a preacher working from that English closer to the
third interpretation. Jerome's and Aquinas's disagreement over that verb, and
Augustine's and Bede's difference over where the weight of the command falls,
therefore do not arise in the speech: it gives Augustine's reach of the
command to each life and Bede's door of each heart, which do not conflict.

The familiar difficulty the speech opens is that the command to watch is
heard every Advent and is known to be badly kept. The textual detail that
changes it is the order of Isaiah 64:6--7: the people confess that none
calls on God's name or rises to take hold of him, and in the next verse they
call on his name, as Father, and name themselves clay. Read with the psalm's
petition that God turn his people and the Collect's petition for the will to
go out to meet Christ, the command to watch becomes something the Church asks
of God and begins to keep by calling on him; Paul's thanksgiving for a
faulty church gives the ground, God's fidelity. The speech states this as
what the texts say and does not stage it as a turn.

Each Father is credited only with what he says of his own passage:
Augustine's reply to a bishop on ``what I say to you, I say to all'', with its
reason; Jerome's patience that only seems the cause of straying, and his
merits and mercy, clay and sons, and ``Be not very angry''; Augustine's one
sentence on the refrain; Bede's authority as the Spirit's grace for good
works, and his door of each heart and account of each; Chrysostom's question
how Paul could praise such a church, his praise that belonged to the grace of
God, his ``God is faithful'' set first against despair, his ``if true'' and his
failure of the call from those who start away; and Augustine's sinner who
confesses, on the Communion verse. The following connections are the
homily's own, or the studies' editorial synthesis, and no Father makes them:
Isaiah 64:6 as the sleep the Gospel warns against (the second
interpretation's inference); the people calling on the name they have just
said none calls on; the psalm beside Isaiah as the prayer of a defeated
people; the Collect's petition for a will set beside the Gospel's command;
Bede's authority as the answer to passivity (the study's comparison);
Isaiah's confession as the pattern of the sacrament of confession, and the
practicable response; the Preface's \latin{vigilantes} as the Gospel's word
(a textual observation the concise study records); the Prayer after
Communion as the forming of the will asked for in the Collect; and the one
ground of Isaiah's people, the Corinthians and the assembly.

Left to the studies: \latin{Descendisti}, Jerome's prayer granted and
Aquinas's three descents; the \work{Ordo}'s titles; the Entrance Antiphon and
its \latin{qui te exspectant}; the psalm's vine and son of man; Augustine on
Christ given as man and to be given as God; Cyril's two comings; Jerome,
Chrysostom and Gregory on why the hour is hidden; Gregory's knock and his
watches of a life; Aquinas on Isaiah 63:17 and on 1~Corinthians 1:7--9;
the inverted order of anger and sin; Jerome's and Aquinas's petitions to turn
and be turned; Bellarmine; the Alleluia verse; the Prayer over the Offerings;
the solemn blessing; and the sacramentary history of the orations.

\textbf{Texts quoted, and their route.} Scripture is quoted in the
Douay--Rheims version as revised by Bishop Challoner, a public-domain study
translation, at Mk 13:37, Is 63:17 and 64:6--7 (Douay 64:7--8), Ps 85
(84):13 and 1~Cor 1:9; the rest is told in the speech's own words. The
English proclaimed in the United States Lectionary is not reproduced; at Is
63:17 it asks why God lets his people wander, so the speech names the Douay
there, whose form of the question is the one Jerome explains. The Collect,
the Preface and the Prayer after Communion are described in part from the
study's account of their Latin, not in the Missal's approved English.
Augustine's reason in the letter to Hesychius is quoted in the study's
English. His sentence on the Communion verse is quoted in the homily's own
English of the Latin below, in which, as in \latin{nisi ille dederit cui
confiteris}, the one confessed to is the one who gives; the studies' English
is not the one spoken. His sentence on the refrain is reported from the
Oxford English of the Nicene and Post-Nicene Fathers, and Chrysostom is
quoted and reported in that series' English, read again in the tracked text
of Homily~2 on First Corinthians. Bede is reported in English that renders
the Latin below, read again in the tracked text; Jerome, from the expansive
study's account of his Latin, which that study read on the Migne page images.
Of Jerome's words only \latin{Ne irascaris, Domine, satis} is given below.

\textbf{Exact loci.} The Gospel's setting and command: Mk 13:3--4, 32,
33--37; the account of the Passion begins at 14:1. Augustine, to whom ``what
I say to you, I say to all'' was said, and \latin{in quo enim quemque
invenerit suus novissimus dies, in hoc eum comprehendet mundi novissimus
dies}: \work{Epistula} 199.1.2--3. The first reading and its prayer: Is
63:16b--17, 19b; 64:4--7 (Douay 63:16b--17; 64:1, 5--8), with 63:18 and
64:10 (Douay 64:11) on the sanctuary and the burnt house, and 64:8 (Douay
64:9), \latin{Ne irascaris, Domine, satis}. Jerome, God's patience and not
God the cause of error and hardness: \work{Commentarii in Isaiam} XVII, PL~24,
cols.~643--644; merits and clemency, clay and sons, and \latin{Ne
irascaris}: cols.~648--650. The psalm: Ps 80 (79):2--3, with the refrain at
vv.~4, 8 and~20. Augustine, ``For averse we have been from Thee, and except
Thou convert us, we shall not be converted'': \work{Enarrationes in Psalmos}
79.3. The Collect: \latin{Da, quaesumus, omnipotens Deus, hanc tuis
fidelibus voluntatem}, with \latin{Christo tuo venienti} and \latin{iustis
operibus occurrentes}.
Bede, \latin{fidelibus suis, concessa sancti Spiritus gratia, facultatem
tribuit bonis operibus serviendi}; \latin{Non apostoli tantum \ldots\ sed et
omnes vigilare praecipimur. Universi januas cordium nostrorum, ne hostis
antiquus mala suggerendo irrumpat, instanter servare jubemur \ldots\
Unusquisque enim pro se reddet rationem Deo}: \work{In Marci evangelium
expositio} IV, on Mk 13:34 and 13:37. The second reading: 1~Cor 1:3--9; the
letter's rebukes at 1:10--12, 3:1, 5:1, 6:1--8 and 11:17--22. Chrysostom,
how Paul praises the Corinthians, ``for not even did this praise belong to
them, but to the grace of God'', ``God is faithful'' put first ``lest after
the vehemence of the reproofs they might fall into despair'', ``\,`God is
faithful,' i.e. `true.' Now if `true,' what things He hath promised He will
also perform'', and ``unbelief cometh not of Him that calleth, but of those
who start away from Him'': \work{Homilies on First Corinthians}~2, on
1~Cor 1:7--9. The Communion Antiphon: Ps 85 (84):13. Augustine, \latin{Tua
ergo peccata potes habere; fructum bonum non habebis, nisi ille dederit cui
confiteris}: \work{Enarrationes in Psalmos} 84.15. Preface I of Advent:
\latin{in humilitate carnis assumptae}; \latin{in suae gloria maiestatis};
\latin{vigilantes nunc audemus exspectare promissum}. The Prayer after
Communion: \latin{Prosint nobis, quaesumus, Domine, frequentata mysteria}.

\par}

% The References stand whole on one page where they can: a list broken after
% one or two entries reads as though the rest had been lost.
\Needspace{16\baselineskip}
\section*{References}

{\small
\begin{itemize}\setlength{\itemsep}{0.15em}\setlength{\parskip}{0pt}
\item \work{Missale Romanum}, editio typica tertia (Typis Vaticanis, 2002), \latin{Dominica I Adventus} and \latin{Praefatio I de Adventu}.
\item \work{Ordo lectionum Missae}, editio typica altera (Libreria Editrice Vaticana, 1981), no.~2, and \work{Lectionary for Mass for Use in the Dioceses of the United States of America}, second typical edition, vol.~I, no.~2, as witnessed by the United States Conference of Catholic Bishops' page of readings for 29~November 2026, read on 26~November 2026.
\item The Holy Bible, Douay--Rheims version, revised by Bishop Richard Challoner (Project Gutenberg text), in its own numbering: Ps 84:13; Is 63:17, 64:7--8; Mk 13:37; 1~Cor 1:9.
\item St Augustine, \work{Enarrationes in Psalmos} 79.3, Nicene and Post-Nicene Fathers, 1st ser., vol.~8; 84.15, Latin text of the Migne printing.
\item St Augustine, \work{Epistula} 199, to Hesychius, Latin text of augustinus.it.
\item St Jerome, \work{Commentarii in Isaiam} XVII, on 63:15--64:9 (Migne, PL 24, Paris, 1865, cols.~643--650).
\item St Bede, \work{In Marci evangelium expositio} IV, on 13:34--37, Latin text of the Migne printing.
\item St John Chrysostom, \work{Homilies on First Corinthians} 2, Nicene and Post-Nicene Fathers, 1st ser., vol.~12.
\end{itemize}
\par}
